\part{Predictive Organization and Interested Action}

\chapter{Predictive Coding, Active Inference, and the Limits of Analogy}

The comparison pursued in this book depends on a technical vocabulary whose terms are often run together, and the first task is to keep them apart. Predictive coding names a family of schemes in which a processing hierarchy transmits predictions downward and discrepancies upward, so that what propagates through the system is chiefly the part of the input that the current prediction fails to account for \citep{rao1999}. Bayesian inference names a normative relation between prior beliefs, a likelihood and a posterior, and says nothing about how a physical system might realize it. Variational inference is an approximation technique for cases in which the exact posterior is intractable, and it proceeds by choosing a tractable family of distributions and selecting the member closest to the posterior. Active inference extends this apparatus by treating action and policy selection as further consequences of the same optimization, so that an agent is described as selecting the courses of action under which its expected observations accord with a generative model that includes preferred outcomes \citep{friston2017,parr2022}. The free-energy principle is the broadest of these claims, a proposal about the conditions under which any self-organizing system that persists can be described as if it were performing such inference \citep{friston2010}. Each step in this sequence adds assumptions, and a conclusion licensed at one step is not licensed at the next. These terms are kept separate throughout, and a comparison with Nietzsche that depends on the stronger claims will say so.

The central mathematical object is the variational free energy. For a generative model $p(o,s)$ over observations $o$ and hidden states $s$, and an approximate posterior $q(s)$, it is defined by
\[
F[q;o]
=
\E_{q(s)}\!\left[\log q(s)-\log p(o,s)\right]
=
\KL\!\left(q(s)\,\|\,p(s\mid o)\right)-\log p(o).
\]
The second equality follows from the factorization $p(o,s)=p(s\mid o)\,p(o)$, which allows the logarithm of the joint to be split into a term that depends on $s$ and a term that does not. Two consequences are immediate. Because the divergence is non-negative, $F$ is an upper bound on the surprise $-\log p(o)$, with equality exactly when $q$ equals the exact posterior. And because $\log p(o)$ does not depend on $q$, minimizing $F$ over $q$ is the same as minimizing the divergence between $q$ and the posterior, so that perception, understood as the choice of $q$, becomes approximate Bayesian inference.

It is equally important to say what does not follow. Minimizing $F$ under a given generative model yields a posterior that is best relative to that model. It carries no guarantee that the model is adequate, that its state space contains the distinctions the situation requires, or that its priors deserve the confidence they receive. The bound on surprise is a bound on surprise as assessed by the model, so that an agent whose model is poor can drive $F$ low by attending to what the model handles well. Approximate inference under a generative model does not by itself establish a philosophical perspectivism, an ontology of power or a criterion of flourishing. Where active inference is invoked in later chapters, it is invoked for what the formalism actually delivers, and the gap between that and the Nietzschean questions about authority and transformation is itself part of the subject.

Action enters through the expected free energy of a policy $\pi$. In one common formulation,
\[
G(\pi)
=
\KL\!\left(q(o\mid\pi)\,\|\,p(o)\right)
+
\E_{q(s\mid\pi)}\!\left[\mathrm{H}\!\left[p(o\mid s)\right]\right],
\]
where the first term measures the divergence between predicted and preferred outcomes, often called risk, and the second measures the expected ambiguity of the mapping from states to observations. Policies with low $G$ are those expected to bring about preferred outcomes through observations that are informative about hidden states, and the decomposition of Appendix~B shows that this informativeness is the expected information gain. The preferred outcomes are encoded in $p(o)$, which is part of the model and not derived from the formalism. This observation is of first importance for what follows, since the authority of those preferences is exactly what the Nietzschean question about value presses on, and Chapter~12 returns to it.

\chapter{Action as the Transformation of Evidence}

An organism does not receive evidence passively. What it encounters depends on where it goes, what it handles, whom it asks and what it ignores, and each of these is in turn shaped by the interpretation it already holds. Perception and action are therefore mutually implicated in a strong sense. Action determines which observations occur, which uncertainties can be resolved and which expectations meet repeated confirmation, and the resulting observations determine how the interpretation is revised. In the formal setting this appears as the dependence of the observation distribution on the policy, written $q(o\mid\pi)$, so that the evidence available for assessing a model is partly a product of the model's own use.

The Nietzschean comparison becomes especially productive at this point. An interpretation can organize practices that make its world increasingly familiar and manageable. A way of reading situations leads to a way of acting in them, the acting produces situations of the kind the reading handles well, and the reading appears to be confirmed. This is the problem of endogenous evidence: apparent confirmation may result partly from restricting encounters to situations that an existing model already handles well. A record of success is evidence about an interpretation only to the extent that the encounters composing it were not selected by the interpretation to ensure that success.

Three kinds of action should be distinguished, because they stand in different relations to the evidential standing of the interpretation they serve. Successful intervention is action that changes a situation in the way the interpretation predicts, and where the situation was not chosen to guarantee that result, it counts as a genuine test that the interpretation passed. Selective exposure is the arrangement of one's circumstances so that only situations of a manageable kind arise, and the confirmations it produces are largely manufactured. Exploratory action seeks situations in which the interpretation's predictions are uncertain or likely to fail, and its results bear on the interpretation in proportion to that risk. The same observable pattern, a long run of confirmed predictions, can arise from any of the three, and nothing in the record of confirmations alone distinguishes them.

The distinction can be given a simple form. Let $\mu_\pi$ be the distribution over situations that an agent following policy $\pi$ actually encounters, and let $\nu$ be a reference distribution over the situations that matter for the agent's continued viability, which may include situations it rarely meets. Write $\ell(x)$ for the model's expected prediction error in situation $x$.

\begin{definition}[Insulation gap]
The \emph{insulation gap} of a policy $\pi$ for a model with error profile $\ell$ is
\[
\Delta(\pi)=\E_{x\sim\nu}\!\left[\ell(x)\right]-\E_{x\sim\mu_\pi}\!\left[\ell(x)\right].
\]
\end{definition}

A large positive gap indicates that the model performs far better on the situations its policy leads it to meet than on those that matter for its viability. This is the signature of selective exposure, but it is equally the signature of effective specialization or of successful environmental management, in which an agent legitimately transforms its surroundings so that the model's competence suffices. The gap measures a disparity in performance relative to a specified reference distribution and does not identify its cause, which requires further evidence about how the encountered situations came to be selected. A gap near zero shows only that the average error is the same under the two distributions. It does not show that the distributions coincide, since very different distributions can yield equal expected values of a single loss. The reference distribution $\nu$ cannot be read off the formalism, and its specification is a substantive and contestable choice. This is the point at which the formal reconstruction acknowledges its dependence on something outside itself, since which situations matter for the agent's continuation is a question about the agent's form of life, and Chapter~17 argues that it must be answered by an account of correction that the formalism does not by itself supply.

The notion also clarifies why defensive coherence is not obviously a failure from within. An agent that has arranged its circumstances so that its interpretation is never seriously challenged will report low error, high confidence and a sense that its view of the world is well supported. Only a comparison with situations it does not meet reveals the gap, and an agent organized to avoid such situations has little occasion to make the comparison.

\chapter{Precision, Attention, and the Authority of an Interpretation}

In predictive models the influence of a discrepancy depends on how it is weighted. A large error that is assigned low precision barely alters the system's estimates, whereas a small error assigned high precision can drive large revisions. Precision is thereby a mechanism of attention and of authority: it determines which discrepancies are heard. In a Nietzschean register the corresponding point is that an encounter does not arrive with an unquestionable authority independent of the organization receiving it. Whether a surprising event counts as a refutation, an anomaly, a nuisance or an opportunity depends on how the interpretation that meets it is disposed to weigh it. Some discrepancies become transformative while others are dismissed, subordinated or redescribed, and the difference lies largely in the weighting.

A local quadratic approximation makes the structure explicit. For a prediction error $\varepsilon$ and a precision matrix $\Pi$, the weighted discrepancy is
\[
\mathcal{E} = \tfrac12\,\varepsilon^{\mathsf T}\Pi\,\varepsilon .
\]
In a Gaussian model with observation noise precision $\Pi$, the negative log-likelihood of an observation has the form
\[
-\log p(o\mid s) = \tfrac12\,\varepsilon^{\mathsf T}\Pi\,\varepsilon-\tfrac12\log\lvert\Pi\rvert+\text{const},
\]
where $\varepsilon$ is the difference between the observation and its prediction. The second term deserves attention. It shows that discounting an error by lowering precision is not free within the formalism, since the log-determinant term penalizes low precision and so charges for any claim that the observations are generally unreliable. A model that could silence every discrepancy by setting $\Pi$ close to zero would be penalized for implying that its own observations carry almost no information. The formalism thus builds in a limited resistance to wholesale insulation, although it cannot prevent a model from raising precision on confirming channels and lowering it on disconfirming ones at a modest total cost.

The philosophical question therefore concerns the formation and revision of $\Pi$ more than the elimination of $\varepsilon$. An account that treats precision as a given parameter to be estimated leaves open who or what assigns authority to some channels of evidence over others, and how that assignment itself becomes open to challenge. If precision is learned from the history of past reliability, then an agent whose history was shaped by selective exposure will have learned precisions that reflect its insulation, and the learning process reproduces rather than corrects the bias of Chapter~6. Revising precision in the face of a discrepancy that the current precision already discounts is thus a distinctive difficulty, and it is one that the literary chapters later examine as a problem of how a text can alter what a reader takes to be authoritative.

This connects the technical notion to reading. Transitions, interruptions, changes of voice and shifts in affect can be studied as ways of altering which aspects of a text a reader treats as authoritative, which is to say as interventions on the reader's precisions. A sudden irony does not supply a new proposition. It changes the weight a reader gives to what came before, and reading afterwards proceeds under a different distribution of attention. The mathematical vocabulary supplies a way of stating this effect, though it does not supply the reasons that make one such shift illuminating and another merely disorienting.

\chapter{Exploration, Resistance, and the Growth of Capacity}

An apparent tension runs between the idea that organisms minimize free energy and the observation that they seek difficulty, novelty and resistance. If the aim were simply to reduce surprise, an organism would retreat to the most predictable environment available, and the Nietzschean praise of resistance would seem opposed to the formalism. The tension is reduced by the expected free energy introduced in Chapter~5. The decomposition of $G(\pi)$ derived in Appendix~B contains the negative expected information gain, so that policies expected to resolve uncertainty about hidden states are rewarded, and exploration is not an exception to the optimization but a component of it. A policy that visits a situation of high uncertainty can have low expected free energy because of what it is expected to teach. Information seeking and the pursuit of unfamiliar situations are therefore accommodated, and the quantity that measures their value is the expected information gain
\[
\mathrm{IG}(\pi)=\E_{q(o\mid\pi)}\!\left[\KL\!\left(q(s\mid o,\pi)\,\|\,q(s\mid\pi)\right)\right],
\]
the expected reduction in uncertainty about hidden states from the observations that policy $\pi$ would produce.

Whether this suffices to reconstruct Nietzschean overcoming is a further question, and the answer depends on a distinction between two kinds of uncertainty. The first is uncertainty resolved within an existing repertoire: the agent has a model with certain variables, does not know their current values, and acts to find out. The second is the encounter that requires a new repertoire, in which the difficulty lies in the fact that the agent's variables omit something essential. Learning a previously unknown value of a familiar variable differs in kind from learning that one's variables were the wrong ones.

The point can be stated sharply. The information gain above is defined over the hidden states $s$ of a fixed model. It is bounded by the entropy of the agent's current beliefs about those states, and it is silent about any feature of the world that the state space does not represent. An unmodelled variable contributes nothing to $\mathrm{IG}$ however much it matters, because the quantity has no argument in which to register it. An agent that explores optimally in the sense of maximizing expected information gain is therefore exploring within a fixed space of distinctions, and its exploration can be exhaustive and still leave the space unchanged.

Nietzsche's interest in resistance can then be investigated as a problem of expanding or reorganizing capacities, and without identifying power with any single information-theoretic quantity. Resistance is valuable on this reading not merely because it supplies information about the familiar variables, though it may, but because it can disclose that the agent's present repertoire cannot meet it, and so create the occasion for a different repertoire. What grows in such an episode is the range of distinctions the agent can draw and the range of situations it can handle, and this is a change in the model family rather than in the estimates made within it. The next part takes up that change directly. Here it is enough to record that active inference handles seeking difficulty as a matter of reducing uncertainty about a given world, and that the Nietzschean demand extends further, to the possibility of a world whose terms differ from those currently available.

A caution against the opposite error belongs here. It would be a mistake to equate overcoming with the maximization of any measure of novelty or difficulty. An agent can pursue difficulty indiscriminately and accomplish nothing, and resistance that exceeds what an organization can absorb destroys it. The growth of capacity is therefore a matter of resistance met at a rate and in a form that the organization can convert into reorganization, and which encounters meet that condition is a question about the particular organism and its supports, which later chapters on attachment and institutions address.
